% ---------------------------------------------------------------------
% 2.9 Redes neuronales artificiales y aprendizaje profundo
% ---------------------------------------------------------------------
\section[REDES NEURONALES ARTIFICIALES Y APRENDIZAJE PROFUNDO]{Redes neuronales artificiales y
aprendizaje profundo}

\subsection{La neurona artificial}

Las redes neuronales artificiales son modelos inspirados en el sistema nervioso, en los que las
unidades de cómputo, llamadas neuronas, se conectan entre sí mediante pesos que cumplen un papel
análogo al de la fuerza de las conexiones sinápticas en los organismos biológicos
\cite{aggarwal2018}. Cada neurona recibe un vector de entradas $\mathbf{x} = (x_1, \dots, x_n)$,
las multiplica por sus pesos $w_i$, suma el resultado junto con un sesgo $b$ y aplica una función
de activación $\varphi$ no lineal \cite{delbrio2002}:
\begin{equation}
y = \varphi\!\left(\sum_{i=1}^{n} w_i x_i + b\right) = \varphi\!\left(\mathbf{w}^{\top}\mathbf{x} + b\right).
\label{eq:neurona}
\end{equation}
La Figura~\ref{fig:neurona} muestra la correspondencia entre este modelo y la neurona biológica
descrita en la Sección~\ref{sec:neuro}: las entradas desempeñan el papel de las dendritas, la suma ponderada
el del soma y la salida el del axón. Al organizar muchas de estas unidades en capas, de manera que
la salida de una capa sea la entrada de la siguiente, se obtiene el perceptrón multicapa, en el que
las capas intermedias se denominan ocultas; al número de capas se le conoce como profundidad de la
red, y de ahí proviene el término aprendizaje profundo \cite{goodfellow2016}.

\begin{figure}[htbp]
\centering
\begin{tikzpicture}[
  ent/.style={circle, draw, minimum size=0.75cm, font=\small},
  flecha/.style={-{Stealth[length=2mm]}, thick}]
  \node[ent] (x1) at (0,1.6) {$x_1$};
  \node[ent] (x2) at (0,0.5) {$x_2$};
  \node at (0,-0.35) {$\vdots$};
  \node[ent] (xn) at (0,-1.2) {$x_n$};
  \node[ent] (b) at (1.9,2.1) {$1$};
  \node[circle, draw, thick, minimum size=1.2cm, fill=gray!10] (sum) at (3.6,0.3) {$\Sigma$};
  \node[rectangle, draw, thick, minimum size=1.0cm, fill=red!8] (act) at (5.6,0.3) {$\varphi(\cdot)$};
  \node (y) at (7.4,0.3) {$y$};
  \draw[flecha] (x1) -- node[above, font=\footnotesize, sloped] {$w_1$} (sum);
  \draw[flecha] (x2) -- node[above, font=\footnotesize, sloped] {$w_2$} (sum);
  \draw[flecha] (xn) -- node[below, font=\footnotesize, sloped] {$w_n$} (sum);
  \draw[flecha] (b) -- node[above right, font=\footnotesize] {$b$} (sum);
  \draw[flecha] (sum) -- node[above, font=\footnotesize] {$z$} (act);
  \draw[flecha] (act) -- (y);
  \node[font=\footnotesize, align=center, text=gray] at (0,-2.2) {Entradas\\(dendritas)};
  \node[font=\footnotesize, align=center, text=gray] at (3.6,-1.25) {Suma ponderada\\(soma)};
  \node[font=\footnotesize, align=center, text=gray] at (5.6,-1.25) {Activación};
  \node[font=\footnotesize, align=center, text=gray] at (7.4,-0.5) {Salida\\(axón)};
\end{tikzpicture}
\caption{Modelo de la neurona artificial y su analogía con la neurona biológica. Elaboración
propia con base en \cite{aggarwal2018}.}
\label{fig:neurona}
\end{figure}

\subsection{Funciones de activación}

Sin la función de activación, una red de cualquier profundidad sería equivalente a una sola
transformación lineal, por lo que esta es la que permite a la red representar relaciones no
lineales \cite{goodfellow2016}. Las funciones sigmoide y tangente hiperbólica fueron las más
utilizadas en las primeras redes, y siguen empleándose en la capa de salida de los clasificadores
binarios y en arquitecturas para P300 como la de Cecotti y Gräser o SepConv1D
\cite{cecotti2011,alvarado2021}; sin embargo, ambas se saturan para valores grandes de su
entrada, lo que hace que sus gradientes tiendan a cero y dificulta el entrenamiento de redes
profundas. La función ReLU resolvió en gran medida este problema y hoy es la opción por defecto en
las capas ocultas \cite{goodfellow2016}, mientras que la ELU, utilizada en EEGNet, permite valores
negativos suaves que acercan la media de las activaciones a cero \cite{lawhern2018}. La
Tabla~\ref{tab:activaciones} presenta su definición y la Figura~\ref{fig:activaciones} su forma.

\begin{table}[htbp]
\centering
\caption{Funciones de activación más utilizadas}
\label{tab:activaciones}
\small
\begin{tabularx}{\textwidth}{L{2.6cm}L{4.3cm}L{1.6cm}Y}
\toprule
\textbf{Función} & \textbf{Definición} & \textbf{Rango} & \textbf{Uso típico} \\
\midrule
Sigmoide & $\sigma(z) = \dfrac{1}{1+e^{-z}}$ & $(0,1)$ & Salida de clasificación binaria (probabilidad de P300). \\[8pt]
Tangente hiperbólica & $\tanh(z) = \dfrac{e^{z}-e^{-z}}{e^{z}+e^{-z}}$ & $(-1,1)$ & Capas ocultas en \cite{cecotti2011,alvarado2021}; redes recurrentes. \\[8pt]
ReLU & $\max(0,z)$ & $[0,\infty)$ & Capas ocultas de redes profundas. \\[4pt]
ELU & $z$ si $z>0$; $\alpha(e^{z}-1)$ si $z\le0$ & $(-\alpha,\infty)$ & Capas ocultas de EEGNet \cite{lawhern2018}. \\[4pt]
Softmax & $\dfrac{e^{z_k}}{\sum_{j} e^{z_j}}$ & $(0,1)$, suma 1 & Salida de clasificación multiclase. \\
\bottomrule
\end{tabularx}
\par\smallskip{\footnotesize Elaboración propia con información de \cite{goodfellow2016,aggarwal2018}.}
\end{table}

\begin{figure}[htbp]
\centering
\includegraphics[width=\textwidth]{02-marco-referencial/funciones_activacion.pdf}
\caption{Gráfica de las funciones de activación sigmoide, tangente hiperbólica, ReLU y ELU.
Elaboración propia.}
\label{fig:activaciones}
\end{figure}

\subsection{Función de pérdida}

El entrenamiento de una red consiste en ajustar sus pesos para minimizar una función de pérdida
que mide la diferencia entre la salida de la red y la etiqueta esperada. Para la clasificación
binaria, con etiquetas $y_i\in\{0,1\}$ y salidas $\hat{y}_i\in(0,1)$, se utiliza la entropía
cruzada binaria \cite{goodfellow2016}:
\begin{equation}
\mathcal{L}_{\text{BCE}} = -\frac{1}{M}\sum_{i=1}^{M}\Big[\,y_i\log\hat{y}_i + (1-y_i)\log(1-\hat{y}_i)\Big].
\label{eq:bce}
\end{equation}
Debido al desbalance de las épocas del P300 Speller, en el que por cada época positiva hay cinco
negativas, es común ponderar el primer término con un factor $\alpha>1$ para que los errores sobre
la clase minoritaria tengan mayor peso, o bien emplear la pérdida focal, que multiplica cada
término por $(1-p_t)^{\gamma}$, donde $p_t$ es la probabilidad asignada a la clase correcta, para
reducir la contribución de los ejemplos que el modelo ya clasifica con facilidad \cite{lin2017}.

\subsection{Descenso de gradiente y retropropagación}

Para minimizar la pérdida se utiliza el método de descenso de gradiente, que actualiza los pesos
en la dirección opuesta al gradiente de la función de pérdida,
\begin{equation}
\mathbf{w}^{(t+1)} = \mathbf{w}^{(t)} - \eta\, \nabla_{\mathbf{w}}\mathcal{L}\big(\mathbf{w}^{(t)}\big),
\label{eq:gradiente}
\end{equation}
donde $\eta$ es la tasa de aprendizaje; si es muy pequeña el entrenamiento resulta lento, y si es
muy grande los pesos oscilan alrededor del mínimo sin alcanzarlo \cite{goodfellow2016}. El cálculo
eficiente de este gradiente en redes de varias capas se realiza mediante el algoritmo de
retropropagación, popularizado por Rumelhart, Hinton y Williams, que propaga el error desde la capa
de salida hacia las capas anteriores aplicando la regla de la cadena \cite{rumelhart1986}. En la
práctica el gradiente se estima sobre pequeños subconjuntos de datos llamados lotes, y en lugar de
una tasa fija se usan optimizadores adaptativos como Adam, que ajusta la tasa de aprendizaje de
cada parámetro a partir de estimaciones del primer y segundo momento de sus gradientes
\cite{kingma2015}, y que fue el utilizado en la comparación de arquitecturas para P300 de
Alvarado-González et al. \cite{alvarado2021}.

\subsection{Regularización}

Dado que las redes profundas tienen una gran cantidad de parámetros y los conjuntos de datos de EEG
son relativamente pequeños, es indispensable emplear técnicas de regularización para evitar el
sobreajuste. Entre las más utilizadas se encuentran el \textit{dropout}, que desactiva
aleatoriamente una fracción de las neuronas en cada iteración de entrenamiento
\cite{srivastava2014}, la normalización por lotes, que normaliza las activaciones de cada capa y
estabiliza el entrenamiento \cite{ioffe2015}, y la parada temprana, que detiene el entrenamiento
cuando la pérdida de validación deja de mejorar \cite{goodfellow2016}. A estas se suma la elección
de una inicialización adecuada de los pesos, como la propuesta por Glorot y Bengio, que mantiene la
varianza de las activaciones similar entre capas \cite{glorot2010}. La Tabla~\ref{tab:regularizacion}
resume estas técnicas junto con su uso en arquitecturas para la detección de la P300.

\begin{table}[htbp]
\centering
\caption{Técnicas de regularización y estabilización del entrenamiento}
\label{tab:regularizacion}
\small
\begin{tabularx}{\textwidth}{L{3.2cm}YY}
\toprule
\textbf{Técnica} & \textbf{Funcionamiento} & \textbf{Uso en P300} \\
\midrule
\textit{Dropout} \cite{srivastava2014} & Anula aleatoriamente una fracción $p$ de las activaciones durante el entrenamiento. & EEGNet y BN3 \cite{lawhern2018,liu2018}. \\
Normalización por lotes \cite{ioffe2015} & Normaliza la media y varianza de las activaciones de cada lote y las reescala con parámetros aprendidos. & BN3, EEGNet \cite{liu2018,lawhern2018}. \\
Regularización $L_2$ & Agrega a la pérdida un término proporcional a $\lVert\mathbf{w}\rVert^2$. & Modelos lineales y redes pequeñas. \\
Parada temprana & Detiene el entrenamiento cuando la pérdida de validación deja de mejorar durante cierto número de épocas. & 50 épocas de paciencia en \cite{alvarado2021}. \\
Inicialización de Glorot \cite{glorot2010} & Inicializa los pesos con una varianza dependiente del número de entradas y salidas de la capa. & Mayoría de las arquitecturas evaluadas en \cite{alvarado2021}. \\
\bottomrule
\end{tabularx}
\par\smallskip{\footnotesize Elaboración propia.}
\end{table}

\subsection{Aprendizaje profundo aplicado al EEG}

El aprendizaje profundo permite que los modelos aprendan representaciones de los datos con
múltiples niveles de abstracción directamente a partir de la señal, sin depender de
características diseñadas a mano \cite{lecun2015}. De acuerdo con Craik et al., un conjunto de
datos de EEG consiste, en su nivel más básico, en una matriz de dos dimensiones, canales por tiempo,
cuya estructura lo hace adecuado para este tipo de técnicas; en su revisión de 90 estudios
encontraron que las redes convolucionales, las recurrentes y las redes de creencia profunda
superaban en precisión a los autocodificadores apilados y a los perceptrones multicapa, y que para
las tareas de ERP se habían utilizado principalmente redes convolucionales, redes de creencia
profunda y autocodificadores \cite{craik2019}. Por su parte, la revisión sistemática de Roy et al.
señala que alrededor del 40\% de los estudios de aprendizaje profundo en EEG emplean redes
convolucionales \cite{roy2019}, mientras que Lotte et al. advierten que, en el momento de su revisión, las redes
profundas aún no mostraban una mejora contundente sobre los mejores métodos clásicos de BCI, en
buena medida por la poca cantidad de datos de entrenamiento con que suele contarse
\cite{lotte2018}.
